\chapter{Related Work}
\label{ch:literature}

\section{Real-Time Object Detection}
\label{sec:lit_objdet}

The detector used in this thesis is YOLO~\cite{redmon2016yolo} as packaged in the Ultralytics framework~\cite{ultralytics2024}. The nano variant was chosen at the start of the project for first-order reasons: real-time throughput on modest hardware, a mature training and deployment ecosystem, and strong small-object performance at the chosen inference resolutions. The cascade then grew around it. It runs an RGB and an IR detector back to back, with a trust classifier and a per-frame confuser filter on top, leaving little room for a heavier backbone. Two-stage architectures such as Faster R-CNN~\cite{ren2015fasterrcnn} and Cascade R-CNN~\cite{cai2018cascade} reach higher accuracy at a latency cost this budget does not allow. The transformer-based detection family, DETR~\cite{carion2020detr} and its real-time variant RT-DETR~\cite{zhao2024rtdetr}, is a plausible drop-in for future work. It was not adopted mid-project for a switching-cost reason that follows from the build order: the downstream stages were trained \emph{after} and \emph{against} the YOLO detectors, so the trust classifier is calibrated to the YOLO confidence-score distribution and the confuser filter reads YOLO's pyramid features. Replacing the backbone now would mean re-deriving both, making it a long-horizon decision rather than a free tuning knob.

A decisive hyperparameter in this work is the inference resolution \texttt{imgsz}. YOLO letterboxes the input to a square canvas; objects smaller than a few pixels on that canvas are unresolvable regardless of model quality, which is why the small-object literature resorts to up-scaled or sliced inference~\cite{akyon2022sahi}. On Svanstr\"om the effect is stark, and it has a concrete geometric cause: the corpus is natively $640\times512$ while the production RGB detector is trained at \texttt{imgsz=1280}, so inference at \texttt{imgsz=640} renders the median 29.8-px drone at half its trained-for scale. A within-model sweep (Section~\ref{sec:resolution_arch}, Figure~\ref{fig:resolution}) quantifies it: the baseline loses $28$~pp of drone recall when inference drops to \texttt{imgsz=640} from the trained-for $1280$, and the aggressively-retrained variant collapses further on the same frames.
% [source: eval/results/svan_resolution_sweep.json; run=eval/svan_resolution_sweep.py] The related flying-object case is treated by \citeauthor{rozantsev2017detecting}~\cite{rozantsev2017detecting}, who detect small UAVs and aircraft from a moving camera by regression-based motion stabilisation of local patches followed by classification on spatio-temporal image cubes, combining appearance with motion cues precisely because single-frame appearance runs out at these object sizes.
% [source: ledger=selcom-imgsz-win; eval=rgb_svan_baseline; run=eval/diagnose_failures.py; cache=eval/results/_failure_diagnosis/svanstrom_1280_by_category.csv; config=svan_iop_1280]
% [source: thesis_eval/results/notes_round1_results.json (svanstrom SZ median 29.8px); appendix: Svanström native 640x512]

\section{Drone Detection and the Confuser Problem}
\label{sec:lit_drone}

At the ranges typical of aerial surveillance, drones, birds, fixed-wing aircraft, and helicopters all present as small, low-texture silhouettes against sky, with similar aspect ratios and little fine detail to separate them. The resulting fire rates are corpus- and category-dependent rather than uniform: on the merged OOD confuser corpus at \texttt{imgsz=640} the baseline detector family fires on 23--58\% of frames per category (Section~\ref{sec:problem}), and the worst case is bird-only Svanstr\"om footage at that benchmark's canonical \texttt{imgsz=1280}, where a YOLO detector trained \emph{with} bird negatives still fires on 94.4\% of frames. This is the \emph{confuser problem}: training-time hard-negative mining suppresses helicopters and partially airplanes but cannot drive the bird fire rate down without sacrificing recall on small-drone targets that share the same silhouette; the one stance that suppresses birds (the aggressively-retrained detector) collapses Svanstr\"om drone recall to $R=0.306$.
% [source: ledger=retrainedv2-recall-collapse; eval=rgb_svan_baseline; thesis_eval/results/notes_round1_results.json (rgb_confuser CAT); cache=eval/results/_failure_diagnosis/svanstrom_1280_by_category.csv; config=svan_iop_1280]

The C-UAS sensing surveys agree that visual detection is operationally attractive (passive, inexpensive, and able to identify what it sees) and that small targets at range and confusion with birds are among its recognised limits~\cite{shi2018counteruas,taha2019drone,samaras2019deep}. \citet{taha2019drone} add an observation this thesis's evaluation protocol responds to directly: published drone-detection results are experimental and ``hardly comparable'' across papers, with reference datasets and requirement-driven problem specifications still missing. Early deep-learning approaches~\cite{aker2017using,rozantsev2017detecting} treated drone detection as a single-class small-object problem; subsequent work explicitly added confuser classes. \citet{schumann2017deep} detect UAVs as a two-stage pipeline (region proposals from background subtraction or a proposal network, then a CNN classifier separating UAVs from distractors such as birds) and show the classifier transfers across a substantial domain gap, an early detect-then-discriminate structure this thesis's cascade echoes. The Drone-vs-Bird challenge~\cite{coluccia2021dronevsbird} consolidated the bird-confuser benchmark; its winning entries relied on explicit negative training images and synthetic data to control bird false positives, and bird confusion remained the challenge's central difficulty across editions.

Two public benchmarks anchor the empirical work. The Svanstr\"om RGB+IR dataset~\cite{svanstrom2021real,svanstrom2022dronedataset} provides per-frame box annotations with three explicit confuser classes (birds, airplanes, helicopters); its native $640\times512$ resolution forces the small-drone regime, making it the discriminating benchmark. The Anti-UAV RGBT benchmark~\cite{jiang2021antiuav,zhao2023antiuav} provides large-scale paired RGB+IR sequences but is saturated for this work's detectors, at both evaluation sizes: the baseline and the aggressively-retrained RGB detector both reach $P=0.9922, R=0.9950, F1=0.9936$ at \texttt{imgsz=1280} ($3{,}178$~TP, $25$~FP, $16$~FN), the production detector reaches $F1=0.986$ at the Tier-1 \texttt{imgsz=640} configuration, and IR-only reaches $F1=0.9654$. The reason is disclosed rather than hidden: Anti-UAV RGB frames appear in both variants' training corpus ($59{,}413$ \texttt{anti}-prefixed frames, Table~\ref{tab:ds_rgb_components}), so no train/test split is claimed; Anti-UAV serves as an in-distribution sanity floor, not a discriminating benchmark. Section~\ref{sec:comparison} compares the production stack to published numbers on both datasets.
% [source: ledger=antiuav-saturated; eval=rgb_antiuav_baseline,v5_antiuav_bare; run=eval/ablate.py; cache=eval/results/_ablation/2026-05-18T22-48-19/master.csv; config=antiuav_iou_1280_s5 + tier1 @640]

The cascade is the response to the residual confuser fire rate that training-time mining cannot eliminate; the trust classifier and per-frame confuser filter catch what the detector cannot.

\section{Thermal Infrared Imaging}
\label{sec:lit_ir}

Thermal infrared cameras image surface emission rather than reflected visible light, with three properties relevant here. First, drone motors, electronic speed controllers (ESCs), and batteries produce localised hot spots against sky, whereas birds tend to maintain more uniform body temperatures, a discriminative cue absent in the visible spectrum and a likely reason the IR detector hallucinates less on bird confusers (Section~\ref{sec:grayscale}). Second, IR is largely illumination-invariant, aiding night and low-visibility operation. Third, single-channel IR imagery shares low-level statistical properties with grayscale RGB: both render small drones as silhouettes against sky. This shared structure underwrites the cross-modal transfer result in Section~\ref{sec:grayscale}: on the YouTube drone-clip diagnostic surface (9 clips, 1{,}359 frames, Appendix~\ref{app:datasets}) the IR detector reaches an aggregate drone $F1=0.636$ on grayscale-converted RGB video, and on its single best clip, \texttt{flock\_of\_seagulls\_attack\_drone\_beach}, it ties the RGB baseline ($0.837$ vs $0.840$); the aggregate, not the best clip, is the representative number. The transfer was emergent, not designed; the grayscale-RGB mode exists as a fallback for when thermal hardware is unavailable.
% [source: ledger=ir-grayscale-fallback; eval=vid_drone_ir_gray (aggregate 0.636; clip tie 0.837/0.840); run=eval/eval_video_tests.py; config=video_drone_iop]

IR imaging also has limitations: lower spatial resolution, thermal contrast collapse when ambient approaches drone surface temperature, and smaller labelled datasets complicating training and OOD evaluation; the IR-side OOD evaluation of this thesis is therefore carried by the thermal confuser corpus and the cross-modal surfaces of Chapter~\ref{ch:hitl}.

\section{Multi-Modal Fusion and Hard-Negative Mining}
\label{sec:lit_fusion}

Multi-modal C-UAS systems combine RGB and thermal streams through early, intermediate, or decision-level fusion~\cite{ramachandram2017fusion}. Intermediate-level fusion (the canonical example being multispectral pedestrian detection~\cite{wagner2016multispectral}) typically requires paired-modality training data. This thesis instead adopts decision-level fusion with a per-frame learned XGBoost trust classifier~\cite{chen2016xgboost,friedman2001greedy}. The motivation is a measured complementarity in \emph{what} the two detectors hallucinate on: the RGB detector's residual fire is bird- and helicopter-driven while the thermal detector's is airplane-dominated, and the thermal channel barely fires on the visible channel's confuser corpus at all. The production router makes its three-way trust decision (RGB, IR, or both) from eight detection-evidence features only (per-modality confidences, box geometry, and cross-modal ratios; Section~\ref{sec:trust_classifier}); a superseded variant that additionally consumed per-frame scene statistics is retained as a comparison (Section~\ref{sec:classifier_compare}). Hard-negative mining at training time (OHEM~\cite{shrivastava2016ohem}, Focal Loss~\cite{lin2017focal}, explicit confuser augmentation) was the first thing tried and is class-asymmetric, as quantified in Section~\ref{sec:problem}: helicopters and partially airplanes respond, birds essentially do not, and the stance that does suppress birds collapses small-drone recall. Birds and small drones are not separable at the detection stage without an unacceptable recall cost; the downstream cascade is the response to that case. The mining experiments reported there are RGB-side by design; the thermal branch's confuser handling is the thermal-native IR filter of Section~\ref{sec:ir_xmodal_verifier}.
% [source: ledger=retrainedv2-recall-collapse; thesis_eval/results/notes_round1_results.json (CAT tables); cache=eval/results/_failure_diagnosis/svanstrom_1280_by_category.csv; config=svan_iop_1280]

\section{Detection Scoring}
\label{sec:lit_calibration}

The choice of matching rule is consequential. Standard IoU under-counts predictions that are correct in centre but imprecise in scale; for small drones whose GT boxes exceed the visible target (Svanstr\"om's annotators drew loose boxes), this is routine, so this thesis uses Intersection over Prediction (IoP) at threshold $0.5$ for Svanstr\"om RGB scoring. The second scoring axis, how detections from two modalities are aggregated into one F1, is defined and justified as evaluation protocol in Section~\ref{sec:scoring_audit}; every multi-modal number in this thesis declares its rule.

\section{Cross-Modal Transfer and Cascaded Rejection}
\label{sec:lit_xmodal_cascade}

Two literatures bear directly on the architectural pieces that do not fit neatly under object detection, thermal imaging, or fusion.

\paragraph{Cross-modal transfer.}
The closest published setting is \emph{modality hallucination}~\cite{hoffman2016modalityhallucination}: training on one modality so the model can be queried at test time on another. RGB-to-thermal style translation~\cite{berg2018rgb2thermal,kniaz2018thermalgan} addresses the inverse direction, and domain-adversarial training~\cite{ganin2016domain} provides the general feature-alignment toolkit. All of these provide or construct cross-modal information at training time. We are not aware of published work matching the setup here, in which no cross-modal information is provided at training time and transfer emerges from shared low-level statistics (single-channel intensity, sky-background silhouettes). The result is presented as emergent in Section~\ref{sec:grayscale}.

\paragraph{Cascaded rejection.}
The architectural pattern of an initial high-recall stage followed by progressively more discriminative rejection predates this application by two decades: \citet{viola2001rapid} for face detection, Cascade R-CNN~\cite{cai2018cascade} for general objects. One framing difference matters: in those designs the first stage is deliberately recall-tuned and weak on precision, whereas the detectors here are strong on both axes standalone (Anti-UAV $P=0.989$/$R=0.982$; 2.8\% hallucination on the composite RGB test corpus) and the cascade exists for a \emph{concentrated} residual failure mode, confuser scenes, that training cannot fix without recall cost. Beyond that, both classical cascades reuse the same backbone across stages, and this design differs in two further ways: (i) the rejection stages are heterogeneous (an XGBoost router on detection-evidence features, then a separate MLP on detection features), making it a refined-decision pipeline, not a refined-bounding-box pipeline; and (ii) the confuser filter can be applied per frame or at the alert-emit boundary (Section~\ref{sec:alert_gate}), and the deployed GUI uses the alert gate. The Viola-Jones intent (early rejection cheap, late rejection expensive) is preserved.

\section{Model-Assisted Annotation and Data-Centric Development}
\label{sec:lit_hitl}

The thesis's human-in-the-loop data engine (Section~\ref{sec:hitl_method}) sits in the active-learning and data-centric-AI lineage. Classical active learning~\cite{settles2009active} selects which samples to label next by model uncertainty; its object-detection instantiations~\cite{brust2019active} rank candidate frames by detection-level criteria. The loop used here inverts the emphasis: rather than selecting unlabelled data for fresh annotation, it routes \emph{model--ground-truth disagreements} to a human adjudicator, treating every disagreement as either a model error (a future hard example) or a label error (a correction). That second branch connects to the label-error literature: confident learning~\cite{northcutt2021confident} shows that benchmark label noise materially distorts model rankings, and the corrections surfaced by this thesis's reviewer (missing boxes, oversized boxes, mis-classed confusers) are exactly the error classes it formalises. The broader data-centric framing~\cite{ng2021datacentric,sambasivan2021data}, that data work rather than model work dominates applied ML outcomes, is borne out quantitatively here: the IR detector's F1 trajectory from $0.503$ to $0.967$ (Section~\ref{sec:ir_evolution}) is driven almost entirely by corpus operations, and its one regression by a corpus operation performed without review.

\section{Probing and Reusing Detector Representations}
\label{sec:lit_probing}

The Model MRI (Section~\ref{sec:model_mri}) belongs to the representation-analysis lineage: linear probes~\cite{alain2016understanding} established reading a frozen network's intermediate features with shallow classifiers as a measurement of what the network already encodes, and feature-space statistics on penultimate layers underpin Mahalanobis-style OOD detection~\cite{lee2018simple}. The MRI applies the same stance to a detection FPN (\texttt{p3}/\texttt{p5} features; classical statistics: LDA, ANOVA, AUROC, and a leakage ratio) but uses the measurements \emph{prescriptively}: where the separation exists, the MRI trains the confuser filter directly on the same detector features (drone detections as positives, confuser detections as negatives). The filter itself is related to knowledge distillation~\cite{hinton2015distilling} in spirit: a small model trained from a larger model's internals. Here, though, the large model is not a teacher producing soft labels but a feature extractor whose representation is reused unchanged, closer to classical transferable-features reuse~\cite{yosinski2014transferable}. Representation-similarity tools such as CKA~\cite{kornblith2019similarity} ask a related question (how similar are two networks' representations); the MRI asks the deployment-facing one: is \emph{this} network's representation sufficient for \emph{that} downstream decision.

\section{Comparison to Prior Work}
\label{sec:comparison}

Direct numerical comparison to other systems is unusually difficult: papers vary in dataset split, inference resolution, scoring rule, and headline metric (mAP, F1, success rate), and this work is not evaluated on the same test splits as any published number. The comparison below is therefore offered as \emph{context, not head-to-head}: a qualitative architectural map (Table~\ref{tab:related_systems}) and a numerical panel for studies reporting metrics on the same public benchmarks (Table~\ref{tab:numerical_comparison}), with the non-comparability caveats stated per row and below the table.

\subsection{Architectural Comparison}

\begin{table}[htbp]
    \centering
    \caption{Architectural map of representative drone-detection systems. ``Modality'' covers what the system consumes; ``Discrimination'' covers how confusers are rejected. The cascade in this thesis is, to our knowledge, the only one that combines learned modality-trust fusion with a downstream per-frame distilled confuser filter.}
    \label{tab:related_systems}
    \footnotesize
    \begin{tabular}{p{2.8cm}p{2.4cm}p{4.2cm}p{3.2cm}}
        \toprule
        System / paper & Modality & Discrimination & Notes \\
        \midrule
        Svanstr\"om 2021~\cite{svanstrom2021real} & RGB + IR (parallel) & Per-modality YOLO only & Dataset paper; no learned fusion \\
        Anti-UAV challenge baselines~\cite{jiang2021antiuav} & RGB + IR (parallel) & Trackers + YOLO; no confuser stage & Drone-tracking task, not raw detection \\
        Coluccia 2021 drone-vs-bird~\cite{coluccia2021dronevsbird} & RGB only & Single-stage YOLO retrained on bird negatives & Hard-negative training, no cascade \\
        Cross-modal attention fusion (e.g., \cite{wagner2016multispectral}) & RGB + IR (intermediate) & End-to-end CNN, no explicit confuser stage & Tightly coupled, requires paired training data \\
        \textbf{This thesis} & RGB + IR (decision-level) & Trust classifier + per-frame distilled (\texttt{mlp\_v5\_v4}) confuser filter & Learned per-frame modality choice; suppression deferred downstream of detection \\
        \bottomrule
    \end{tabular}
\end{table}

Prior work uses either a single detector or a tightly-coupled multi-modal one and treats confuser discrimination as a training-time problem. This thesis keeps the detectors strong on both precision and recall in their own right and adds a separate discrimination stage (trust classifier $+$ confuser filter, run per frame or at the alert gate) for the concentrated confuser failure mode, making the modality-trust decision per frame.

\subsection{Numerical Comparison}

\begin{table}[htbp]
    \centering
    \caption{Numerical comparison on public benchmarks. \textbf{All comparisons require the caveats below the table.} ``Scoring'' indicates the matching rule used in the cited paper where it can be determined; where it cannot, the cell is marked \texttt{ND} (not disclosed). The Svanstr\"om $P$/$R$ are means of the paper's per-distance-bin values (it reports per-bin $P$/$R$ and an average $F1$), so they do not exactly satisfy $F1=2PR/(P+R)$.}
    \label{tab:numerical_comparison}
    \footnotesize
    \resizebox{\textwidth}{!}{%
    \begin{tabular}{p{2.7cm}p{3.6cm}p{3.0cm}ccc}
        \toprule
        Dataset & System & Scoring & $P$ & $R$ & F1 \\
        \midrule
        \multirow{2}{2.6cm}{Svanstr\"om (visible video)} & Svanstr\"om 2021 (avg.\ over distance bins) & IoU@0.5, det.\ thr.\ 0.5, YOLOv2 @416 & 0.837 & 0.748 & 0.785 \\
                                                   & \textbf{This thesis (baseline RGB, S1)} & IoP@0.5, @1280 & 0.940 & 0.961 & 0.950 \\
        \midrule
        \multirow{2}{2.6cm}{Svanstr\"om (thermal video)}  & Svanstr\"om 2021 (avg.\ over distance bins) & IoU@0.5, det.\ thr.\ 0.5, YOLOv2 @256 & 0.843 & 0.661 & 0.760 \\
                                                   & \textbf{This thesis (IR \texttt{v3b}, \texttt{imgsz=640})} & IoP@0.5 & 0.950 & 0.973 & 0.961 \\
        \midrule
        \multirow{2}{2.6cm}{Anti-UAV (drone, RGB)}    & Anti-UAV baselines (range) & IoU & ND & ND & 0.90--0.97 \\
                                                   & \textbf{This thesis (\texttt{ft4}, @640)} & IoU@0.5 & 0.989 & 0.982 & 0.985 \\
        \midrule
        \multirow{2}{2.6cm}{Anti-UAV (drone, IR)}     & Anti-UAV baselines (range) & IoU & ND & ND & 0.90--0.96 \\
                                                   & \textbf{This thesis (IR \texttt{v3b}, @640)} & IoU@0.5 & 0.966 & 0.956 & 0.961 \\
        \bottomrule
    \end{tabular}}
\end{table}
% [source: thesis_eval/results/tier1_results.json (ir_confusers C: bare 0.2943, clf->filt[robust6] 0.0192 with thermal-native cbam filter); honest held-out IR_confusers val/test shipped 90 -> 22, see Section~\ref{sec:ir_xmodal_verifier}]
% [source: svanstrom2021real (arXiv 2007.07396): visible avg F1 0.7849 (close 0.8682), IR avg F1 0.7601 (close 0.8845); YOLOv2, inputs 416x416/256x256, bboxPrecisionRecall IoU 0.5, det thr 0.5; eval set = 120 IR + 120 visible clips (5 per class x distance bin), 240 clips train]
% [source: svanstrom2021real per-bin P/R verified via ar5iv 2026-06: Vcam P 0.8989/0.8391/0.7726 R 0.7355/0.7306/0.7785; IRcam P 0.9197/0.8282/0.7822 R 0.8737/0.7039/0.4043; filled table P/R = mean over distance bins (Vcam 0.837/0.748, IRcam 0.843/0.661); F1 column = paper avg = mean of per-bin F1]

Three caveats apply to every row of Table~\ref{tab:numerical_comparison}.

\emph{Scoring rule.} Svanstr\"om \etal report per-sensor precision/recall via MATLAB's \texttt{bboxPrecisionRecall} (an IoU-overlap matcher) at IoU $0.5$ and a detection threshold of $0.5$, with YOLOv2 at input sizes $416{\times}416$ (visible) and $256{\times}256$ (thermal); their table values are averages over distance bins, and their close-range bins are higher ($0.868$ visible, $0.885$ IR). The Anti-UAV baselines rarely disclose their IoU threshold, and neither source discloses a multi-modal aggregation rule. IoP@0.5, used here for Svanstr\"om RGB, is more lenient than IoU@0.5; an IoU@0.5 re-score would produce lower numbers.

\emph{Dataset split.} The Svanstr\"om paper held out an evaluation set of 120 IR and 120 visible clips (five per class and distance bin), but the clip list is not published with the dataset, so it cannot be adopted here. This thesis instead evaluates the full Svanstr\"om corpus, which is leakage-free on the RGB side (no Svanstr\"om frames of any class appear in any RGB training corpus of this thesis, Section~\ref{sec:svanstrom_audit}) and disclosed as in-distribution on the IR side. One further conversion caveat: the dataset's annotations ship as MATLAB \texttt{groundTruth} objects, and this thesis converted them to YOLO format with its own script, so box geometry may differ marginally from the authors' MATLAB pipeline. Anti-UAV has a more canonical split but uses sequence-level tracking metrics not directly equivalent to per-frame F1.

\emph{Task.} Anti-UAV baselines are tracking systems. Their state-accuracy numbers (0.90--0.97 range across years) include temporal aggregation absent from this thesis's frame-level numbers, and no per-frame precision or recall is reported, so those cells stay \texttt{ND}. The Anti-UAV row is included for orientation only.
% [source: ledger=antiuav-saturated; eval=rgb_antiuav_baseline; run=eval/ablate.py; cache=eval/results/_ablation/2026-05-18T22-48-19/master.csv; config=antiuav_iou_1280_s5]
% [source: github.com/DroneDetectionThesis/Drone-detection-dataset (no published split; MATLAB groundTruth labels)]

Two claims survive these caveats. First, this thesis's per-modality numbers \emph{exceed} the published Svanstr\"om averages ($0.950$ vs $0.785$ visible, $0.961$ vs $0.760$ IR); the caveats above carry the qualification, since part of that gap is attributable to detector generation (YOLO26n vs YOLOv2), inference resolution ($1280$/$640$ vs $416$/$256$), and matcher leniency (IoP vs IoU) rather than to a like-for-like superiority. At minimum, the detectors are competitive baselines rather than weak strawmen. Second, no published system reports a separately-evaluated OOD confuser fire rate at the granularity of Table~\ref{tab:ablation_confusers}, because no published system explicitly separates a confuser-rejection stage from drone detection. The suppression numbers of Section~\ref{sec:pipeline_confusers} ($30.4\% \to 4.9\%$ at the router, $\to 1.4\%$ at the per-frame filter, at the canonical configuration) therefore cannot be directly benchmarked. The contribution is the architectural separation that makes such numbers measurable.
% [source: thesis_eval/results/tier1_results.json (rgb_confuser C_confuser)]

% ======================================================================
% CHAPTER 3 — METHODOLOGY
% ======================================================================
